\documentclass[11pt,reqno]{amsart}
\usepackage[T1]{fontenc}
\usepackage{amsmath,amssymb,amsthm,mathtools,booktabs}
\usepackage[margin=1in]{geometry}
\usepackage[colorlinks=true,linkcolor=blue,citecolor=blue,urlcolor=blue]{hyperref}
\newtheorem{theorem}{Theorem}[section]
\newtheorem{proposition}[theorem]{Proposition}
\newtheorem{lemma}[theorem]{Lemma}
\newtheorem{corollary}[theorem]{Corollary}
\theoremstyle{definition}
\newtheorem{definition}[theorem]{Definition}
\newtheorem{example}[theorem]{Example}
\newtheorem{exercise}[theorem]{Exercise}
\theoremstyle{remark}
\newtheorem{remark}[theorem]{Remark}
\newcommand{\R}{\mathbb R}
\newcommand{\F}{\mathcal F}
\newcommand{\NN}{\mathcal N}
\newcommand{\Id}{I}
\DeclareMathOperator{\diag}{diag}
\DeclareMathOperator{\arsinh}{arsinh}
\title[Finite Spin to the Harmonic Oscillator]{From Finite Spin to the Harmonic Oscillator:\\Krawtchouk Transforms, Mehler Kernels, and Convolution}
\subjclass[2020]{33C45, 22E70, 47D06, 60J60, 42A38}
\keywords{Krawtchouk polynomials, Hermite functions, harmonic oscillator, Mehler kernel, $\mathfrak{sl}_2$, contraction, Ornstein--Uhlenbeck process, Hamming scheme}
\date{Lecture notes}
\begin{document}
\begin{abstract}
These lecture notes develop the passage from the finite-dimensional
(spin) representations of $\mathfrak{sl}_2$ to the quantum harmonic
oscillator. On the finite side the protagonists are the Krawtchouk
polynomials, the Krawtchouk transform, and the Ehrenfest urn chain; on
the continuum side they are the Hermite functions, the Fourier transform,
and the Ornstein--Uhlenbeck semigroup, whose kernel is Mehler's. We
treat the oscillator semigroup in real and imaginary time, general
Gaussian kernels as squeezed oscillator semigroups, the contraction of
spin to the Heisenberg algebra, the scaling of Krawtchouk to Hermite
polynomials, and the fate of Krawtchouk convolution in the limit.
Exercises are included at the end of each lecture.
\end{abstract}
\maketitle
\tableofcontents

\section*{Introduction}
The finite-dimensional irreducible representations of $\mathfrak{sl}_2$
have a distinguished family of orthogonal polynomials attached to them:
the Krawtchouk polynomials, introduced as discrete analogues of the
Hermite polynomials \cite{Krawtchouk1929}. In matrix form they appear as
group elements of the spin representations \cite{FeinsilverKocik2005},
and in coding theory as the eigenvalues of the Hamming association scheme
\cite{Delsarte1973,MacWilliamsSloane1977}. The harmonic oscillator, by
contrast, lives on the infinite-dimensional space $L^2(\R)$, where the
Hermite functions diagonalize simultaneously the Hamiltonian and the
Fourier transform. The aim of these notes is to make precise the many
ways in which the finite picture tends to the infinite one, and to be
equally precise about where the analogy fails.

The overall plan may be summarized as follows.
\[
\begin{array}{ccccc}
\text{finite spin}&\longrightarrow&\text{Heisenberg algebra}
 &\longrightarrow&\text{quadratic (metaplectic) $\mathfrak{sl}_2$}\\
\text{Krawtchouk polynomials}&\longrightarrow&\text{Hermite polynomials}
 &\longrightarrow&\text{Hermite functions}\\
\text{Ehrenfest chain}&\longrightarrow&\text{Ornstein--Uhlenbeck}
 &\longrightarrow&\text{oscillator heat semigroup}\\
\text{finite bilinear kernels}&\longrightarrow&\text{OU transition kernel}
 &\longrightarrow&\text{Mehler kernel}.
\end{array}
\]
The first arrow in each row is a limit as $N\to\infty$; the second is
an exact construction (forming quadratic expressions, multiplying by a
Gaussian, or conjugating by a unitary map).

Lectures~\ref{L:hermite}--\ref{L:gauss} are purely about the continuum.
Lectures~\ref{L:spin}--\ref{L:limits} set up the finite theory and pass to
the limit. Lecture~\ref{L:conv} discusses convolution, and
Lecture~\ref{L:end} treats the degenerate Gaussian kernels at the edge
of the oscillator region.

\subsection*{Standing conventions}
$D=d/dx$. Operators on $L^2(\R)$ are first defined on the Schwartz
space $\mathcal S(\R)$ and then closed; $H$ always denotes the
self-adjoint oscillator realization. All square roots of positive numbers
are positive. The Fourier transform is
\begin{equation}\label{eq:F}
(\F f)(\xi)=\frac1{\sqrt{2\pi}}\int_\R e^{-ix\xi}f(x)\,dx,\qquad
(\F^*f)(\xi)=\frac1{\sqrt{2\pi}}\int_\R e^{+ix\xi}f(x)\,dx,
\end{equation}
so $\F^*=\F^{-1}$. Both are unitary on $L^2(\R)$.

%======================================================================
\section{Hermite functions, the oscillator, and the Fourier transform}\label{L:hermite}

\subsection{Two-parameter Hermite polynomials}
\begin{definition}
For $\tau\in\R$ define polynomials $H_n(x,\tau)$ by
\begin{equation}\label{eq:Hgen}
e^{zx-\tau z^2/2}=\sum_{n\ge0}H_n(x,\tau)\frac{z^n}{n!}.
\end{equation}
Write $H_n(x)=H_n(x,1)$ for the probabilists' Hermite polynomials.
\end{definition}
Replacing $z$ by $z\sqrt\tau$ shows $H_n(x,\tau)=\tau^{n/2}H_n(x/\sqrt\tau)$
for $\tau>0$; in particular $H_n(x,\tfrac12)=2^{-n/2}H_n(\sqrt2x)$.
Differentiating \eqref{eq:Hgen} with respect to $x$, to $z$, and to
$\tau$ gives
\begin{equation}\label{eq:Hrec}
\partial_xH_n=nH_{n-1},\qquad
H_{n+1}=xH_n-n\tau H_{n-1},\qquad
\partial_\tau H_n+\tfrac12\partial_x^2H_n=0.
\end{equation}
Combining the first two, $\tau H_n''-xH_n'+nH_n=0$.

\subsection{Hermite functions and ladder operators}
Put
\[
\psi_n(x)=2^nH_n(x,\tfrac12)e^{-x^2/2},\qquad
H=\tfrac12(x^2-D^2).
\]
\begin{proposition}\label{prop:ladder}
For all $n\ge0$, with $\psi_{-1}=0$,
\[
(x+D)\psi_n=2n\psi_{n-1},\qquad
(x-D)\psi_n=\psi_{n+1},\qquad
H\psi_n=(n+\tfrac12)\psi_n.
\]
Consequently $\psi_n=(x-D)^ne^{-x^2/2}$ and
\begin{equation}\label{eq:psigen}
\sum_{n\ge0}\psi_n(x)\frac{u^n}{n!}=e^{-x^2/2+2ux-u^2}.
\end{equation}
\end{proposition}
\begin{proof}
Write $H_n=H_n(x,\frac12)$. By the product rule
$(x+D)(f e^{-x^2/2})=f'e^{-x^2/2}$ and
$(x-D)(fe^{-x^2/2})=(2xf-f')e^{-x^2/2}$. The first and second
identities follow from \eqref{eq:Hrec} with $\tau=\frac12$, since
$2xH_n-nH_{n-1}=2H_{n+1}$. Then
$2H=(x+D)(x-D)-1$ gives $2H\psi_n=(2n+1)\psi_n$.
The generating function is \eqref{eq:Hgen} with $z=2u$, $\tau=\frac12$.
\end{proof}

The normalized Hermite functions
\[
h_n=\frac{\psi_n}{\pi^{1/4}\sqrt{2^nn!}}
=\frac{H_n(\sqrt2x)e^{-x^2/2}}{\pi^{1/4}\sqrt{n!}}
\]
form an orthonormal basis of $L^2(\R,dx)$. With
\[
a=\frac{x+D}{\sqrt2},\qquad a^\dagger=\frac{x-D}{\sqrt2},\qquad
\NN=a^\dagger a=H-\tfrac12\Id,
\]
one has $[a,a^\dagger]=\Id$, $ah_n=\sqrt n\,h_{n-1}$,
$a^\dagger h_n=\sqrt{n+1}\,h_{n+1}$, and $\NN h_n=nh_n$.
The self-adjoint operator $H$ is defined spectrally by this basis.

\subsection{The Fourier transform as a quarter period}
\begin{theorem}\label{thm:fourier}
$\F h_n=(-i)^nh_n$ and $\F^*h_n=i^nh_n$. Consequently
\begin{equation}\label{eq:Fosc}
\F=e^{-i\pi\NN/2},\qquad
e^{-i\pi H/2}=e^{-i\pi/4}\F,\qquad
e^{i\pi H/2}=e^{i\pi/4}\F^*.
\end{equation}
\end{theorem}
\begin{proof}
The Gaussian integral
$(2\pi)^{-1/2}\int e^{zx-x^2/2}dx=e^{z^2/2}$, valid for complex $z$,
applied to \eqref{eq:psigen} gives
\[
\F^*\Big(\sum\psi_n\tfrac{u^n}{n!}\Big)(\xi)
=e^{-u^2}e^{(i\xi+2u)^2/2}=e^{-\xi^2/2+2i\xi u+u^2}
=\sum\psi_n(\xi)\frac{(iu)^n}{n!}.
\]
Comparing coefficients yields $\F^*\psi_n=i^n\psi_n$, and taking
inverses $\F\psi_n=(-i)^n\psi_n$. Both sides of \eqref{eq:Fosc} are
unitary and agree on an orthonormal basis.
\end{proof}

\begin{remark}[Signs]
The sign in \eqref{eq:F} matters. If the positive-exponent transform is
called $\F$, then $\F$ and $\F^*$ must be exchanged everywhere, and the
formula reads $e^{i\pi H/2}=\sqrt i\,\F$ with $\sqrt i=e^{i\pi/4}$.
In particular, with \eqref{eq:F}, the correct statement is
$e^{i\pi H/2}=\sqrt i\,\F^{-1}$. A common source of error is a lost
factor $2$ in the mixed term: the exponent is $2i\xi u$, not $i\xi u$.
\end{remark}
More generally $e^{-i\theta\NN}$ is the fractional Fourier transform of
angle $\theta$. Its discrete counterpart built from Krawtchouk functions
is studied in \cite{AtakishievWolf1997}; we meet it in
Lecture~\ref{L:limits}.

\subsection{The quadratic $\mathfrak{sl}_2$}
On $\mathcal S(\R)$ put
\begin{equation}\label{eq:quad}
\mathsf R=\tfrac12x^2,\qquad \mathsf L=\tfrac12D^2,\qquad \rho=xD+\tfrac12.
\end{equation}
Using $[D,x]=1$ one checks
\begin{equation}\label{eq:quadbr}
[\mathsf L,\mathsf R]=\rho,\qquad[\mathsf L,\rho]=2\mathsf L,\qquad
[\rho,\mathsf R]=2\mathsf R,\qquad H=\mathsf R-\mathsf L.
\end{equation}
The closure of $\rho$ is skew-adjoint and generates the unitary
dilations $(e^{s\rho}f)(x)=e^{s/2}f(e^sx)$. Thus
$\{\mathsf L,\mathsf R,\rho\}$ generate rotations (Fourier),
dilations (squeezing), and, in imaginary time, heat semigroups. These
quadratic operators preserve parity, so this representation splits into
even and odd pieces; it is very different from a finite irreducible
representation, a point we return to in Lecture~\ref{L:limits}.

\begin{exercise}
Verify \eqref{eq:quadbr}, and show that
$K_+=\frac12(a^\dagger)^2$, $K_-=\frac12a^2$, $K_0=\frac12H$ satisfy
$[K_0,K_\pm]=\pm K_\pm$, $[K_-,K_+]=2K_0$. Find the lowest $K_0$
weights on even and odd functions.
\end{exercise}
\begin{exercise}
Show that $\F x\F^{-1}=-D_\xi/i$ in the sense
$\widehat{xf}=iD_\xi\widehat f$ and $\widehat{Df}=i\xi\widehat f$, and
deduce $\F(x-D)\F^{-1}=-i(x-D)$. Use this to give a second proof of
Theorem~\ref{thm:fourier}.
\end{exercise}

%======================================================================
\section{The oscillator semigroup and Mehler's formula}\label{L:mehler}

\subsection{The heat kernel of $H$}
\begin{theorem}\label{thm:heat}
For $s>0$, $e^{-sH}$ is the integral operator on $L^2(\R,dx)$ with kernel
\begin{equation}\label{eq:heat}
K_s(x,y)=\frac1{\sqrt{2\pi\sinh s}}
\exp\left[-\frac{(x^2+y^2)\cosh s-2xy}{2\sinh s}\right].
\end{equation}
\end{theorem}
\begin{proof}[Derivation]
The kernel must satisfy $\partial_sK=\frac12(\partial_x^2-x^2)K$ with
$K_s\to\delta(x-y)$ as $s\downarrow0$. The symmetric Gaussian ansatz
$K=A(s)\exp[-\frac a2(x^2+y^2)+bxy]$ gives
$a'=1-a^2=-b^2$, $b'=-ab$, $A'/A=-a/2$, and the initial condition
forces $a\sim b\sim1/s$, $A\sim(2\pi s)^{-1/2}$. Thus
$a=\coth s$, $b=1/\sinh s$, $A=(2\pi\sinh s)^{-1/2}$.
A complete proof that this kernel represents $e^{-sH}$ follows from
Corollary~\ref{cor:mehlerL2} below, which verifies the action on the
Hermite basis.
\end{proof}

\begin{proposition}\label{prop:onbasis}
For $t<0$, with $K_{-t}$ as in \eqref{eq:heat},
\[
\int_\R K_{-t}(x,y)\psi_n(y)\,dy=e^{t(n+1/2)}\psi_n(x).
\]
\end{proposition}
\begin{proof}
Integrate the kernel against the generating function \eqref{eq:psigen}.
A Gaussian integral gives
$e^{t/2}e^{-x^2/2+2e^tux-e^{2t}u^2}$, which is
$\sum_n\frac{u^n}{n!}e^{t(n+1/2)}\psi_n(x)$. Compare coefficients.
\end{proof}

\subsection{Positive time and analytic continuation}
For real $t>0$ the spectral operator $e^{tH}$ is unbounded, with domain
\[
\mathcal D(e^{tH})=\Big\{\sum c_nh_n:\ \sum e^{2t(n+1/2)}|c_n|^2<\infty\Big\},
\]
which contains every $h_n$, and $e^{tH}h_n=e^{t(n+1/2)}h_n$. However, the
kernel \eqref{eq:heat} with $s=-t<0$ does \emph{not} give a convergent
integral, even on $\psi_n$: after multiplying by $\psi_n(y)$ the
coefficient of $y^2$ in the exponent is
$\frac12(\coth t-1)=(e^{2t}-1)^{-1}>0$. For complex $t$ with
$\operatorname{Re}t<0$, the kernel
\[
\frac1{\sqrt{-2\pi\sinh t}}
\exp\left[\frac{(x^2+y^2)\cosh t-2xy}{2\sinh t}\right]
\]
is analytic, with the root continued from the positive root on
$t<0$. Its boundary value at $t=i\pi/2$ is
$(2\pi)^{-1/2}e^{i\pi/4}e^{ixy}$, recovering \eqref{eq:Fosc}.
Boundary values are understood as oscillatory integrals or strong
limits, not as absolutely convergent kernels on all of $L^2$.

\subsection{Mehler's formula}
\begin{theorem}[Mehler \cite{Mehler1866}]\label{thm:mehler}
For real $x,y$ and complex $|r|<1$,
\begin{equation}\label{eq:mehler}
\sum_{n\ge0}\frac{r^n}{n!}H_n(x)H_n(y)
=\frac1{\sqrt{1-r^2}}
\exp\left(\frac{2rxy-r^2(x^2+y^2)}{2(1-r^2)}\right),
\end{equation}
the root being the branch equal to $1$ at $r=0$. The series converges
absolutely and locally uniformly in the disk, and its radius of
convergence is exactly $1$.
\end{theorem}
\begin{proof}
For real $0\le r<1$ and Gaussian measure
$d\gamma(z)=(2\pi)^{-1/2}e^{-z^2/2}dz$, use
$\int e^{uz-u^2/2}e^{vz-v^2/2}d\gamma(z)=e^{uv}$, i.e.\
$\int H_mH_n\,d\gamma=n!\,\delta_{mn}$, together with
$H_n(x)=\int(x+iz)^n\,d\gamma(z)$. Then
\[
\sum_n\frac{r^n}{n!}H_n(x)H_n(y)
=\iint\exp\big(r(x+iz)(y+iw)\big)\,d\gamma(z)\,d\gamma(w),
\]
which is a Gaussian integral equal to the right side of \eqref{eq:mehler}.
Both sides are analytic in $|r|<1$, so the identity extends.
At $r=\pm1$ the right side is singular, so the radius is $1$.
\end{proof}

\begin{corollary}\label{cor:mehlerL2}
For $0<r<1$,
\begin{equation}\label{eq:mehlernorm}
\sum_{n\ge0}r^nh_n(x)h_n(y)
=\frac1{\sqrt{\pi(1-r^2)}}
\exp\left[-\frac{1+r^2}{2(1-r^2)}(x^2+y^2)+\frac{2r}{1-r^2}xy\right],
\end{equation}
and this is the kernel of $r^{\NN}$ on $L^2(\R,dx)$. With $r=e^{-s}$,
multiplying by $e^{-s/2}$ gives \eqref{eq:heat}.
\end{corollary}
\begin{proof}
Substitute $h_n$ into \eqref{eq:mehler} at $(\sqrt2x,\sqrt2y)$.
The kernel is square-integrable, hence defines a bounded operator,
which agrees with $r^{\NN}$ on the Hermite basis.
\end{proof}

\begin{remark}[Which operator is Mehler's sum?]
The bare sum \eqref{eq:mehler} is a kernel relative to $\gamma$:
\[
(T_rf)(x)=\int M_r(x,y)f(y)\,d\gamma(y)
\]
satisfies $T_rH_n=r^nH_n$, so $T_r=r^{L}$ with $L=-D^2+xD$ on
$L^2(\gamma)$. This is the Ornstein--Uhlenbeck semigroup. The unitary
$f\mapsto(2\pi)^{-1/4}e^{-x^2/4}f(x)$ followed by the dilation
$x\mapsto\sqrt2x$ intertwines $L$ with $\NN$. Thus Mehler's sum is
$r^{H-1/2}$ \emph{after} Gaussian conjugation and rescaling.
\end{remark}

\begin{exercise}
Show $\|e^{tH}\|=e^{t/2}$ for $t<0$, and that $e^{tH}$ is trace class
with $\operatorname{tr}e^{tH}=1/(2\sinh(-t/2))$. Verify this by
integrating \eqref{eq:heat} along the diagonal.
\end{exercise}
\begin{exercise}
Show that the semigroup property $K_s*K_{s'}=K_{s+s'}$ (composition of
kernels) reduces to the addition formulas for $\cosh$ and $\sinh$.
\end{exercise}

%======================================================================
\section{Gaussian kernels as squeezed oscillator semigroups}\label{L:gauss}
Consider, for $a>0$ and real $b$, the operator $A_{a,b}$ with kernel
\[
G_{a,b}(x,y)=e^{-a(x^2+y^2)+bxy}.
\]
\begin{theorem}\label{thm:gauss}
Suppose $|b|<2a$ and set
\[
\omega=\sqrt{4a^2-b^2},\qquad r=\frac{b}{2a+\omega},\qquad
C=\sqrt{\frac{\pi(1-r^2)}{\omega}}.
\]
Let $H_\omega=\frac12(-D^2+\omega^2x^2)$,
$\NN_\omega=H_\omega/\omega-\frac12$, and
$h^{(\omega)}_n(x)=\omega^{1/4}h_n(\sqrt\omega x)$. Then
\[
G_{a,b}(x,y)=C\sum_{n\ge0}r^nh^{(\omega)}_n(x)h^{(\omega)}_n(y).
\]
Thus $A_{a,b}$ is trace class with eigenvalues $Cr^n$;
$A_{a,b}=Cr^{\NN_\omega}$ for $r>0$,
$A_{a,b}=CP|r|^{\NN_\omega}$ for $r<0$ (with $Pf(x)=f(-x)$), and
$A_{a,0}$ is rank one.
\end{theorem}
\begin{proof}
One checks $a=\frac{\omega(1+r^2)}{2(1-r^2)}$ and
$b=\frac{2\omega r}{1-r^2}$, and then applies \eqref{eq:mehlernorm} in
the variables $\sqrt\omega x,\sqrt\omega y$. For $r<0$ use
$Ph_n=(-1)^nh_n$.
\end{proof}
No translation is needed, because the quadratic form has no linear
terms. The rescaling is the unitary dilation
$V_\omega=e^{\frac12(\log\omega)\rho}$, with $\NN_\omega=V_\omega\NN V_\omega^{-1}$.
So $A_{a,b}$ is a function of a \emph{squeezed} oscillator, not of $H$
itself. For $0<b<2a$, writing $r=e^{-\tau}$,
\begin{equation}\label{eq:gaussheat}
a=\frac\omega2\coth\tau,\qquad b=\frac\omega{\sinh\tau},\qquad
A_{a,b}=\sqrt{\frac{2\pi\sinh\tau}{\omega}}\,e^{-\tau H_\omega/\omega}.
\end{equation}
In the Lie-algebraic form of \eqref{eq:quad},
$e^{-\tau H_\omega/\omega}=\exp[\tau(\omega^{-1}\mathsf L-\omega\mathsf R)]$.

\begin{example}\label{ex:sqrt8}
For $a=\frac32$, $b=1$ one finds $\omega=\sqrt8$, $r=3-\sqrt8$,
$\tau=2\log(1+\sqrt2)$. Kernels of exactly this type arise as candidate
central limits in the spectral analysis of the symmetric Krawtchouk
matrices of Lecture~\ref{L:spin}, whose spectra were studied in
\cite{FeinsilverFitzgerald1996}. The theorem identifies the Gaussian
$e^{-\sqrt8x^2/2}$ as the \emph{exact} ground state of such a limiting
operator. Transferring this to the finite matrices is a separate,
genuinely asymptotic problem: compact-window kernel asymptotics give
lower bounds via Rayleigh quotients, but convergence of the full spectrum
requires global estimates.
\end{example}

\begin{exercise}
Show that if $|b|>2a$ the kernel does not define a bounded operator on
$L^2(\R)$.
\end{exercise}

%======================================================================
\section{Finite spin and Krawtchouk polynomials}\label{L:spin}

\subsection{The spin representations}
Fix $N\ge1$ and a basis $e_0,\dots,e_N$. Define
\[
R_Ne_j=(N-j)e_{j+1},\qquad L_Ne_j=je_{j-1},\qquad
\Lambda_Ne_j=(N-2j)e_j,
\]
with out-of-range terms zero. Then
\begin{equation}\label{eq:finbr}
[L_N,R_N]=\Lambda_N,\qquad [\Lambda_N,R_N]=-2R_N,\qquad
[\Lambda_N,L_N]=2L_N.
\end{equation}
This is the irreducible representation of dimension $N+1$; on
homogeneous polynomials $e_j=\xi^{N-j}\eta^j$ it is
$R=\eta\partial_\xi$, $L=\xi\partial_\eta$,
$\Lambda=\xi\partial_\xi-\eta\partial_\eta$. Comparing with
\eqref{eq:quadbr}, the dictionary
\begin{equation}\label{eq:dict}
\mathsf L\leftrightarrow-L_N,\qquad\mathsf R\leftrightarrow R_N,\qquad
\rho\leftrightarrow-\Lambda_N
\end{equation}
matches the bracket relations exactly. This is an isomorphism of
abstract Lie algebras only: the finite and quadratic representations
have different adjoint structures and are not unitarily equivalent.

With $B_N=\diag\binom Nj$, conjugation by $B_N^{1/2}$ gives orthonormal
coordinates in which
\[
\widehat R_N|j\rangle=\sqrt{(N-j)(j+1)}\,|j+1\rangle,\qquad
\widehat L_N|j\rangle=\sqrt{j(N-j+1)}\,|j-1\rangle,
\]
and $\widehat L_N=\widehat R_N^*$. Thus $J_+=\widehat L_N$,
$J_-=\widehat R_N$, $J_z=\frac12\Lambda_N$ is the spin-$\frac N2$
representation, with $m=\frac N2-j$.

\subsection{Krawtchouk polynomials}
\begin{definition}
The (symmetric) Krawtchouk polynomials $K_n(k;N)$ are defined by
\begin{equation}\label{eq:Kgen}
(1+z)^{N-k}(1-z)^k=\sum_{n=0}^NK_n(k;N)z^n,\qquad 0\le k\le N.
\end{equation}
\end{definition}
Thus $K_0=1$, $K_1=N-2k$, $K_2=\frac12((N-2k)^2-N)$. With the binomial
weight $\pi_N(k)=2^{-N}\binom Nk$,
\begin{equation}\label{eq:Korth}
\sum_{k=0}^N\pi_N(k)K_n(k;N)K_m(k;N)=\binom Nn\delta_{mn},
\end{equation}
and the duality $\binom NkK_n(k;N)=\binom NnK_k(n;N)$ holds.
Write $p_{n,N}=K_n/\sqrt{\binom Nn}$ for the orthonormal polynomials.

\subsection{The Krawtchouk transform}
Let $(\Phi_N)_{nk}=K_n(k;N)$. From \eqref{eq:Kgen} one finds
$\Phi_N^2=2^N\Id$; equivalently $\Phi_N$ is the matrix of the
substitution $(\xi,\eta)\mapsto(\xi+\eta,\xi-\eta)$ in the spin
representation \cite{FeinsilverKocik2005}. The orthogonal, symmetric
version is
\begin{equation}\label{eq:Q}
Q_N=2^{-N/2}B_N^{-1/2}\Phi_NB_N^{1/2},\qquad
(Q_N)_{kn}=\sqrt{\pi_N(k)}\,p_{n,N}(k),\qquad Q_N^2=\Id.
\end{equation}
The \emph{symmetric Krawtchouk matrix} $S_N=\Phi_NB_N$ is symmetric but
not orthogonal; its spectrum is a different and harder problem
\cite{FeinsilverFitzgerald1996}.

\begin{proposition}\label{prop:Qdiag}
Let $M_N=\frac N2\Id-\frac12(\widehat R_N+\widehat L_N)$. Then
$M_NQ_N=Q_N\diag(0,1,\dots,N)$.
\end{proposition}
\begin{proof}
This is the three-term recurrence of the $K_n$ in the variable $k$,
rewritten in orthonormal coordinates (equivalently, the difference
equation of Lecture~\ref{L:limits}).
\end{proof}
Thus $\widehat R_N+\widehat L_N=2J_x$ plays the role of a position
operator, and $M_N$ is a ``finite oscillator'' Hamiltonian with exactly
equally spaced spectrum $0,1,\dots,N$; compare
\cite{AtakishievSuslov1990,AtakishievWolf1997}.

\begin{exercise}
Check $\Phi_2^2=4\Id$ and $Q_2^2=\Id$ explicitly, and verify
Proposition~\ref{prop:Qdiag} for $N=2$.
\end{exercise}
\begin{exercise}
Show that the Casimir $J_x^2+J_y^2+J_z^2$ equals
$\frac N2(\frac N2+1)\Id$ and explain why no finite-dimensional
representation can satisfy $[a,a^\dagger]=\Id$.
\end{exercise}

%======================================================================
\section{From spin to oscillator: contraction and diffusion}\label{L:limits}

\subsection{Contraction to the Heisenberg algebra}
Set $a_N=\widehat L_N/\sqrt N$, $a_N^\dagger=\widehat R_N/\sqrt N$,
$\NN_N=\frac12(N\Id-\Lambda_N)$. Then
\begin{equation}\label{eq:contr}
[a_N,a_N^\dagger]=\Id-\frac{2}{N}\NN_N,\qquad
[\NN_N,a_N]=-a_N,\qquad[\NN_N,a_N^\dagger]=a_N^\dagger.
\end{equation}
For fixed $j$, $a_N|j\rangle\to\sqrt j\,|j-1\rangle$ and
$a_N^\dagger|j\rangle\to\sqrt{j+1}\,|j+1\rangle$. Under $|j\rangle\mapsto h_j$
this is the oscillator representation of Lecture~\ref{L:hermite}. The
convergence holds on each fixed finite set of levels, near the
highest-weight end; it is not uniform over all levels, as the trace
obstruction to $[a,a^\dagger]=\Id$ demands.

The quadratic algebra \eqref{eq:quad} is then obtained \emph{after} the
limit, as quadratic expressions in $a,a^\dagger$. The spin algebra does
not contract to the quadratic $\mathfrak{sl}_2$; it contracts to the
Heisenberg algebra, from which $\mathfrak{sl}_2$ is rebuilt.

\subsection{Krawtchouk to Hermite}
Use the central coordinate and mesh
\begin{equation}\label{eq:scale}
x_k=\frac{N-2k}{\sqrt{2N}},\qquad\delta_N=\sqrt{2/N}.
\end{equation}
\begin{theorem}\label{thm:KH}
If $x_{k_N}\to x$, then for each fixed $n$,
\[
n!\,(2N)^{-n/2}K_n(k_N;N)\to H_n(x,\tfrac12),\qquad
p_{n,N}(k_N)\to\frac{H_n(\sqrt2x)}{\sqrt{n!}},
\]
and $\delta_N^{-1/2}(Q_N)_{k_Nn}\to h_n(x)$.
\end{theorem}
\begin{proof}
Put $z=u/\sqrt{2N}$ in \eqref{eq:Kgen}. With $N-k=\frac N2+\sqrt{N/2}\,x_k$
and $k=\frac N2-\sqrt{N/2}\,x_k$,
\[
(N-k)\log(1+z)+k\log(1-z)=xu-\tfrac{u^2}4+o(1),
\]
locally uniformly in complex $u$. Cauchy estimates give coefficientwise
convergence to \eqref{eq:Hgen} with $\tau=\frac12$. The de Moivre--Laplace
local limit $\pi_N(k_N)/\delta_N\to\pi^{-1/2}e^{-x^2}$ gives the last
statement.
\end{proof}
Limit relations of this type between discrete and continuous classical
families are discussed systematically in \cite{GodoyEtAl1998}.

\subsection{Ehrenfest to Ornstein--Uhlenbeck}
The Ehrenfest urn chain has generator
\begin{equation}\label{eq:GN}
(G_Nf)(k)=\frac{N-k}2\big(f(k+1)-f(k)\big)+\frac k2\big(f(k-1)-f(k)\big).
\end{equation}
It is reversible for $\pi_N$, and $G_Np_{n,N}=-np_{n,N}$. Equivalently,
each of $N$ independent bits flips at rate $\frac12$.
\begin{theorem}\label{thm:OU}
For smooth $F$ and $f(k)=F(x_k)$,
$G_Nf(k)\to(GF)(x)$ where $G=\frac12D^2-xD$. The invariant measures
converge to $d\mu=\pi^{-1/2}e^{-x^2}dx$. The unitary
$U:L^2(\mu)\to L^2(dx)$, $Uf=\pi^{-1/4}e^{-x^2/2}f$, satisfies
\[
-UGU^{-1}=H-\tfrac12\Id=\NN.
\]
At finite $N$, $U_Nf(k)=\sqrt{\pi_N(k)}f(k)$ gives $-U_NG_NU_N^{-1}=M_N$.
\end{theorem}
\begin{proof}
Taylor expansion with step $\delta_N$: the drift is
$\frac{N-2k}{2}\cdot(-\delta_N)\to-x$ and the diffusion coefficient is
$\frac N2\cdot\frac{\delta_N^2}{2}\to\frac12$. The conjugation is a
direct computation: $e^{x^2/2}(\frac12D^2-xD)e^{-x^2/2}=\frac12D^2-\frac12x^2+\frac12$.
\end{proof}
Note that the Hermite polynomials appear in two different stochastic
roles. In the fixed-$N$ urn above they are eigenfunctions of a
\emph{stationary} diffusion. Along a simple random walk $X_m$, by
contrast, the polynomials $K_n$ with order parameter $m$ give discrete
martingales, whose limits $H_n(B_\tau,\tau)$ are Brownian martingales by
the heat equation in \eqref{eq:Hrec}. These are distinct scalings.

\subsection{Finite fractional transforms}
Define $\mathfrak F_{N,\theta}=Q_N\diag(e^{-in\theta})Q_N^T=e^{-i\theta M_N}$.
These are unitary, and on each fixed limiting mode they converge to
$e^{-i\theta\NN}$; at $\theta=\frac\pi2$ the limit is $\F$. This is the
fractional Fourier--Krawtchouk transform of \cite{AtakishievWolf1997}.
Note that $Q_N$ itself is an involution while $\F$ has order four;
$Q_N$ is the analogue of a reflection in the spin group, not of the
Fourier transform.

\begin{exercise}
Show that with the opposite coordinate $x_k=(2k-N)/\sqrt{2N}$ every
limit in Theorem~\ref{thm:KH} acquires a factor $(-1)^n$.
\end{exercise}

%======================================================================
\section{Finite bilinear kernels and the Mehler limit}\label{L:kernels}
We recall the notation of Lecture~\ref{L:spin}. The binomial weight is
\[
\pi_N(\ell)=2^{-N}\binom N\ell,\qquad 0\le\ell\le N,
\]
and the orthonormal Krawtchouk polynomials are
\[
p_{n,N}(k)=\binom Nn^{-1/2}K_n(k;N),\qquad
(1+z)^{N-k}(1-z)^k=\sum_{n=0}^NK_n(k;N)z^n,
\]
so that $\sum_k\pi_N(k)p_{n,N}(k)p_{m,N}(k)=\delta_{mn}$ by
\eqref{eq:Korth}. For $0<r<1$, set
\begin{equation}\label{eq:finmehler}
m_{r,N}(k,\ell)=\sum_{n=0}^Nr^np_{n,N}(k)p_{n,N}(\ell)
=\sum_{n=0}^N\frac{r^n}{\binom Nn}K_n(k;N)K_n(\ell;N),
\qquad 0\le k,\ell\le N.
\end{equation}
Explicitly (see Proposition~\ref{prop:finmehler}),
\[
m_{r,N}(k,\ell)=\frac{2^N}{\binom N\ell}\sum_{j}
\binom kj\binom{N-k}{\ell-j}
\Big(\frac{1+r}2\Big)^{N-k-\ell+2j}
\Big(\frac{1-r}2\Big)^{k+\ell-2j},
\]
the sum running over $\max(0,k+\ell-N)\le j\le\min(k,\ell)$. Here $j$ is
the number of the $k$ initially occupied sites that remain occupied.
\begin{proposition}\label{prop:finmehler}
With $r=e^{-s}$, $P_s(k,\ell)=\pi_N(\ell)m_{r,N}(k,\ell)$ is the
transition matrix $e^{sG_N}$. It is nonnegative and
\[
\sum_\ell P_s(k,\ell)z^\ell
=\Big(\tfrac{1+r}2+\tfrac{1-r}2z\Big)^{N-k}
 \Big(\tfrac{1-r}2+\tfrac{1+r}2z\Big)^{k}.
\]
Moreover $\sum_j\pi_N(j)m_{r,N}(k,j)m_{q,N}(j,\ell)=m_{rq,N}(k,\ell)$.
\end{proposition}
\begin{proof}
The spectral expansion follows from $G_Np_{n,N}=-np_{n,N}$. In the bit
picture, each bit is flipped at time $s$ with probability $\frac{1-r}2$,
which gives the generating function. The composition law is
orthonormality \eqref{eq:Korth}.
\end{proof}
Nonnegativity of such kernels is the finite counterpart of the
positivity of Mehler's kernel; compare the characterization of positive
definite Krawtchouk sequences in \cite{Eagleson1969}.

\begin{theorem}\label{thm:finlimit}
If $x_{k_N}\to x$, the law of $x_{\ell}$ under $P_s(k_N,\cdot)$
converges to that of $rx+\sqrt{1-r^2}\,Z$ with $Z\sim\mu$. Equivalently,
the limiting transition operator is
\[
(T_rf)(x)=\int f\big(rx+\sqrt{1-r^2}z\big)\,d\mu(z),
\]
whose density relative to $\mu$ is
\[
m_r(x,y)=\frac1{\sqrt{1-r^2}}
\exp\left(\frac{2rxy-r^2(x^2+y^2)}{1-r^2}\right),
\]
and $Um_rU^{-1}$ is the kernel \eqref{eq:mehlernorm} of $r^{\NN}$.
\end{theorem}
\begin{proof}
The generating function in Proposition~\ref{prop:finmehler} is a product
of independent Bernoulli factors; the central limit theorem gives the
stated Gaussian with mean $rx$ and variance $\frac{1-r^2}2$. This avoids
interchanging the $N\to\infty$ limit with the infinite spectral sum.
\end{proof}
So Mehler's formula is the $N\to\infty$ shadow of the elementary
statement that independent bit flips form a Markov semigroup.

\begin{exercise}
Derive \eqref{eq:mehlernorm} from Theorem~\ref{thm:finlimit} and the
fixed-mode limits of Theorem~\ref{thm:KH}, justifying the interchange
of limit and sum by a domination argument.
\end{exercise}

%======================================================================
\section{Convolution: what survives the limit}\label{L:conv}

\subsection{Krawtchouk convolution}
Identify $\{0,1\}^N$ with subsets of $\{1,\dots,N\}$; the group law is
symmetric difference. If independent uniform subsets have sizes
$k,\ell$, the size of their sum is $k+\ell-2J$, with $J$ hypergeometric.
This defines the radial (hypergroup) convolution
\begin{equation}\label{eq:hyper}
\delta_k*_N\delta_\ell=\sum_j
\frac{\binom kj\binom{N-k}{\ell-j}}{\binom N\ell}\,\delta_{k+\ell-2j}.
\end{equation}
Its identity is $\delta_0$. The normalized characters
$\chi_{n,N}(k)=K_n(k;N)/\binom Nn$ satisfy the product formula
\begin{equation}\label{eq:char}
\int\chi_{n,N}\,d(\delta_k*_N\delta_\ell)=\chi_{n,N}(k)\chi_{n,N}(\ell),
\end{equation}
so the transform $\widehat\nu(n)=\sum_k\chi_{n,N}(k)\nu(k)$ turns $*_N$
into multiplication. These are the eigenvalues of the Hamming
association scheme \cite{Delsarte1973,MacWilliamsSloane1977}; the
product formula is a special case of the Krawtchouk addition theorem
\cite{Dunkl1976}.

\subsection{Central inputs: collapse}
\begin{proposition}\label{prop:collapse}
If $x_{k_N}\to x$ and $x_{\ell_N}\to y$, then the pushforward of
$\delta_{k_N}*_N\delta_{\ell_N}$ to the $x$-coordinate converges to $\mu$,
independently of $x,y$.
\end{proposition}
\begin{proof}
The mean of $x_M$ is $\sqrt{2/N}\,x_{k_N}x_{\ell_N}\to0$ and its variance
tends to $\frac12$; apply the hypergeometric central limit theorem.
\end{proof}
Two central inputs thus produce the invariant Gaussian, forgetting both.
Spectrally, $\chi_{n,N}(k_N)=O(N^{-n/2})$, and the factor $N^{n/2}$ needed
to obtain Hermite polynomials appears twice on the right of
\eqref{eq:char} but once on the left. The reason is geometric: the
identity $k=0$ has moved to $x_0=\sqrt{N/2}\to\infty$.

\subsection{One macroscopic input: the Mehler kernel}
\begin{theorem}\label{thm:mixed}
Let $x_{k_N}\to x$ and $\ell_N/N\to\frac{1-r}2$ with $-1<r<1$. Then
\[
x_M\Longrightarrow rx+\sqrt{1-r^2}\,Z,\qquad Z\sim\mu,
\]
and $\chi_{n,N}(\ell_N)\to r^n$ for each fixed $n$.
\end{theorem}
Now one factor in \eqref{eq:char} stays of order one, and the product
formula has a nontrivial limit: the Ornstein--Uhlenbeck kernel. There
is an exact finite version. The radial binomial measure
\[
\nu_{r,N}(\ell)=\binom N\ell\Big(\frac{1-r}2\Big)^\ell\Big(\frac{1+r}2\Big)^{N-\ell}
\]
(flip each bit with probability $\frac{1-r}{2}$) has character
multipliers $r^n$, so convolution by $\nu_{r,N}$ \emph{is} the finite
Mehler operator of Proposition~\ref{prop:finmehler}, and
$\nu_{r,N}*_N\nu_{q,N}=\nu_{rq,N}$. In the limit this becomes
\[
\int m_r(x,z)m_q(z,y)\,d\mu(z)=m_{rq}(x,y),
\]
i.e.\ the semigroup law $e^{-sH}e^{-tH}=e^{-(s+t)H}$.

\subsection{Other products}
Ordinary convolution on $\R$ is addition of independent variables,
diagonalized by characters $e^{ix\xi}$. The OU operation instead
contracts one input and adds independent noise. The Hermite polynomials
satisfy the binomial identity
\[
H_n(ax+by,\tfrac12)=\sum_j\binom nja^jb^{n-j}H_j(x,\tfrac12)H_{n-j}(y,\tfrac12),
\qquad a^2+b^2=1,
\]
which splits degrees rather than multiplying characters. On polynomials
one can also define the Wick product
$H_m(\cdot,\frac12)\diamond H_n(\cdot,\frac12)=H_{m+n}(\cdot,\frac12)$,
equivalently $f\diamond g=W((W^{-1}f)(W^{-1}g))$ with $W=e^{-D^2/4}$.
It is an algebraic product, not a positive convolution.

\begin{remark}[An intertwining question]
One may ask for an intertwiner between ordinary convolution with the
Fourier transform and Krawtchouk convolution with the Krawtchouk
transform. The discussion above shows that the question needs
qualification: a finite convolution theorem exists, and a natural
Mehler/OU limit exists, but under central scaling Krawtchouk convolution
does not tend to ordinary convolution. Any proposed intertwiner must
specify the product, the measure, and the topology to be preserved.
\end{remark}

\begin{exercise}
Compute the mean and variance of $k+\ell-2J$ and prove
Proposition~\ref{prop:collapse} and Theorem~\ref{thm:mixed}.
\end{exercise}
\begin{exercise}
Verify that $\nu_{r,N}$ has character multipliers $r^n$, using
\eqref{eq:Kgen} with $z$ replaced by a suitable ratio.
\end{exercise}

%======================================================================
\section{The degenerate Gaussian kernels: $b=\pm2a$}\label{L:end}
We return to Lecture~\ref{L:gauss} at the boundary of the region
$|b|<2a$, where the oscillator picture breaks down.

\subsection{$b=2a$: free heat flow}
Here $G_{a,2a}(x,y)=e^{-a(x-y)^2}$, so $A_{a,2a}$ is ordinary convolution
with a Gaussian. By \eqref{eq:F},
$\F(A_{a,2a}f)(\xi)=\sqrt{\pi/a}\,e^{-\xi^2/(4a)}\F f(\xi)$, hence
\[
A_{a,2a}=\sqrt{\pi/a}\;e^{D^2/(4a)}.
\]
It is bounded with norm $\sqrt{\pi/a}$, but not compact; its spectrum is
$[0,\sqrt{\pi/a}]$, purely continuous. Ordinary convolution appears
exactly when the confining quadratic form degenerates.

\subsection{$b=-2a$: heat flow with reflection}
Here $G_{a,-2a}(x,y)=e^{-a(x+y)^2}$ and
\[
A_{a,-2a}=\sqrt{\pi/a}\;e^{D^2/(4a)}P,
\]
with $P$ commuting with the heat operator. It equals the previous
operator on even functions and its negative on odd functions; its
spectrum is $[-\sqrt{\pi/a},\sqrt{\pi/a}]$.

\subsection{Approach to the endpoints}
As $b\uparrow2a$ with $a$ fixed, $\omega\to0$, $r\to1$, and
$\tau/\omega\to1/(2a)$. In \eqref{eq:gaussheat} the potential
$\frac{\omega^2x^2}{2}$ disappears, the kinetic part tends to
$D^2/(4a)$, and the scalar tends to $\sqrt{\pi/a}$. The uniform Schur
bound
\[
\sup_x\int e^{-a(x^2+y^2)+bxy}\,dy=\sqrt{\frac\pi a}\,
\sup_xe^{-(4a^2-b^2)x^2/(4a)}=\sqrt{\frac\pi a}\qquad(|b|\le2a)
\]
and pointwise convergence on compactly supported functions give strong
convergence $A_{a,b}\to A_{a,2a}$ on $L^2$. Norm convergence fails,
since the interior operators are compact and the endpoint operator is
not. The case $b\downarrow-2a$ follows by composing with $P$. For
$|b|>2a$ the form is indefinite, and the kernel no longer defines a
bounded operator.

\begin{exercise}
Show that the eigenvalues $Cr^n$ of $A_{a,b}$ accumulate densely in
$[0,\sqrt{\pi/a}]$ as $b\uparrow2a$, consistent with the spectrum of
the endpoint operator.
\end{exercise}

%======================================================================
\begin{thebibliography}{99}
\bibitem{AtakishievSuslov1990}
N.~M. Atakishiev and S.~K. Suslov, \emph{Difference analogs of the
harmonic oscillator}, Theoret. and Math. Phys. \textbf{85} (1990).
\href{https://doi.org/10.1007/BF01017247}{doi:10.1007/BF01017247}.

\bibitem{AtakishievWolf1997}
N.~M. Atakishiyev and K.~B. Wolf, \emph{Fractional Fourier--Kravchuk
transform}, J. Opt. Soc. Amer. A \textbf{14} (1997), no.~7, 1467.
\href{https://doi.org/10.1364/JOSAA.14.001467}{doi:10.1364/JOSAA.14.001467}.

\bibitem{Delsarte1973}
P.~Delsarte, \emph{An algebraic approach to the association schemes of
coding theory}, Philips Research Reports Supplements, No.~10 (1973).

\bibitem{Dunkl1976}
C.~F. Dunkl, \emph{A Krawtchouk polynomial addition theorem and wreath
products of symmetric groups}, Indiana Univ. Math. J. \textbf{25}
(1976), 335--358.

\bibitem{Eagleson1969}
G.~K. Eagleson, \emph{A characterization theorem for positive definite
sequences of the Krawtchouk polynomials}, Australian J. Statist.
\textbf{11} (1969), 29--38.

\bibitem{FeinsilverFitzgerald1996}
P.~Feinsilver and R.~Fitzgerald, \emph{The spectrum of symmetric
Krawtchouk matrices}, Linear Algebra Appl. \textbf{235} (1996), 121--139.
\href{https://doi.org/10.1016/0024-3795(94)00123-5}{doi:10.1016/0024-3795(94)00123-5}.

\bibitem{FeinsilverKocik2005}
P.~Feinsilver and J.~Kocik, \emph{Krawtchouk polynomials and Krawtchouk
matrices}, in: \emph{Recent Advances in Applied Probability}
(R.~Baeza-Yates, J.~Glaz, H.~Gzyl, J.~H\"usler, and J.~L.~Palacios, eds.),
Springer, New York, 2005, pp.~115--141.
\href{https://doi.org/10.1007/0-387-23394-6_5}{doi:10.1007/0-387-23394-6\_5}.

\bibitem{GodoyEtAl1998}
E.~Godoy, A.~Ronveaux, A.~Zarzo, and I.~Area, \emph{On the limit
relations between classical continuous and discrete orthogonal
polynomials}, J. Comput. Appl. Math. \textbf{91} (1998), 97--105.
\href{https://doi.org/10.1016/S0377-0427(98)00026-0}{doi:10.1016/S0377-0427(98)00026-0}.

\bibitem{Krawtchouk1929}
M.~Krawtchouk, \emph{Sur une g\'en\'eralisation des polyn\^omes
d'Hermite}, C. R. Acad. Sci. Paris \textbf{189} (1929), 620--622.

\bibitem{MacWilliamsSloane1977}
F.~J. MacWilliams and N.~J.~A. Sloane, \emph{The Theory of
Error-Correcting Codes}, North-Holland, Amsterdam, 1977.

\bibitem{Mehler1866}
F.~G. Mehler, \emph{Ueber die Entwicklung einer Function von beliebig
vielen Variablen nach Laplaceschen Functionen h\"oherer Ordnung},
J. Reine Angew. Math. \textbf{66} (1866), 161--176.
\href{https://doi.org/10.1515/crll.1866.66.161}{doi:10.1515/crll.1866.66.161}.
\end{thebibliography}
\end{document}
